% ============================================================
%  Abstract  (~200 words, no section heading per LNCS style)
% ============================================================
Cross-framework comparisons of quantum software development kits (SDKs)
are methodologically compromised by a simulator confound: each framework
is typically evaluated on its own native backend, making it impossible
to attribute observed differences to the transpiler rather than the
simulator.
We resolve this by introducing a unified benchmarking pipeline built on
QCanvas, which compiles circuits from any source framework to OpenQASM
3.0 and routes all execution through a single QSim statevector backend.
We apply this pipeline to 12 quantum algorithms across 3 frameworks
(Qiskit, Cirq, PennyLane), up to 8 qubits, at 4096 shots, and across
6 metric dimensions.
Three findings stand out.
First, Qiskit and Cirq are virtually equivalent in compiled circuit
quality (QPack $S_{\text{overall}}$ 58.80 vs.\ 58.79), differing only
in ancilla-placement strategy on error-correction circuits.
Second, PennyLane exhibits super-linear gate scaling on oracle and
search algorithms: for Grover's algorithm the power-law exponent reaches
$a{=}2.21$, yielding 216 gates at 6 qubits versus 18 for Qiskit and
Cirq.
Third, Jensen--Shannon divergence equals exactly 1.000 between PennyLane
and both other frameworks for six algorithms at 4 or more qubits,
indicating fully non-overlapping measurement distributions---a
systematic decomposition incompatibility, not noise.
NISQ practitioners should treat Qiskit and Cirq as interchangeable in
circuit-quality terms and select PennyLane only when its differentiable
programming model is required, accepting the overhead at scale.

\keywords{quantum benchmarking \and framework comparison \and
OpenQASM 3.0 \and NISQ \and circuit transpilation}
